\chapter{Conclusion Générale}
\label{chap:conclusion}

\section{Bilan des objectifs}

Au cours de ces cinq semaines de stage chez Axeli, j'ai eu l'opportunité de concevoir et de mettre en oeuvre un laboratoire de sécurité réseau complet utilisant des pare-feux FortiGate dans un environnement virtualisé GNS3.

\begin{table}[H]
\centering
\caption{Bilan final des objectifs}
\label{tab:bilan}
\begin{tabularx}{\textwidth}{|c|X|c|}
\hline
\textbf{N°} & \textbf{Objectif} & \textbf{Résultat} \\
\hline
1 & firewall policies & Réalisé \\
2 & NAT (SNAT) & Réalisé \\
3 & Routage statique & Réalisé \\
4 & Routage inter-LAN & Réalisé \\
5 & Routage BGP & Non réalisé \\
6 & VPN IPSec & Non réalisable \\
7 & VPN SSL & Non réalisable \\
8 & High Availability & Réalisé \\
9 & Antivirus & Partiel (signatures factory) \\
10 & IPS & Partiel (signatures factory) \\
11 & Application Control & Non réalisable \\
12 & Web Filtering & Partiel (listes statiques) \\
13 & DNS Filtering & Partiel (listes statiques) \\
14 & SSL Inspection & Partiel (certificat généré, non appliqué) \\
15 & Authentification locale & Réalisé \\
16 & Logging & Réalisé \\
\hline
\end{tabularx}
\end{table}

Sur les 16 objectifs, 7 ont été entièrement réalisés, 5 partiellement (profils configurés mais non testables sans FortiGuard) et 4 ne sont pas réalisables (contraintes de licence ou d'infrastructure).

\section{Compétences acquises}

\subsection{Compétences techniques}

\begin{itemize}
    \item Configuration et administration de pare-feux FortiGate (FortiOS 7.4)
    \item Mise en place de routage dynamique (OSPF, BGP partiel)
    \item Déploiement de profils de sécurité (AV, IPS, Web Filter, SSL Inspection)
    \item Virtualisation réseau avec GNS3 et Docker
    \item Administration de bases de données PostgreSQL
    \item Développement de serveurs d'application (Flask)
    \item Mise en place de solutions de supervision (Grafana, Prometheus)
\end{itemize}

\subsection{Compétences méthodologiques}

\begin{itemize}
    \item Conception et documentation de topologies réseau
    \item Résolution de problèmes techniques complexes
    \item Travail en environnement virtualisé
    \item Rédaction de documentation technique
\end{itemize}

\subsection{Compétences professionnelles}

\begin{itemize}
    \item Intégration en entreprise (Axeli)
    \item Travail en équipe avec des ingénieurs certifiés
    \item Respect des procédures et des bonnes pratiques
    \item Présentation de résultats techniques
\end{itemize}

\section{Difficultés rencontrées}

Les principales difficultés que j'ai rencontrées sont les suivantes :

\begin{enumerate}
    \item \textbf{Contraintes de licence} : les limites de l'évaluation ont nécessité des solutions de contournement ingénieuses
    \item \textbf{Gestion des VMs} : les performances limitées (1 vCPU, 2 GB RAM) ont impacté le temps de configuration
    \item \textbf{Docker dans GNS3} : la chaîne d'initialisation (init.sh -- run-cmd.sh) a nécessité des ajustements
    \item \textbf{Compromis HA/OSPF} : le choix entre HA et OSPF a été contraint par le nombre d'interfaces
    \item \textbf{VPN cloud} : l'interconnexion IPSec avec un endpoint cloud n'a pas pu être finalisée en raison des restrictions de licence et des contraintes du firewall virtuel
\end{enumerate}

\section{Perspectives}

\subsection{Améliorations possibles}

\begin{itemize}
    \item Obtention de licences FortiGuard pour des fonctionnalités complètes
    \item Ajout d'un FortiAnalyzer pour la supervision centralisée
    \item Déploiement de FortiManager pour la gestion centralisée
    \item Extension du laboratoire avec des équipements supplémentaires
\end{itemize}

\subsection{Déploiement en entreprise}

Je suis convaincu que les connaissances acquises lors de ce stage sont directement applicables dans un contexte professionnel :
\begin{itemize}
    \item Déploiement de pare-feux FortiGate en production
    \item Configuration de politiques de sécurité et NAT
    \item Mise en place de solutions de supervision
    \item Supervision et réponse aux incidents
\end{itemize}

\subsection{Projet professionnel}

Ce stage a confirmé mon intérêt pour les métiers de la sécurité réseau. Les compétences acquises en configuration FortiGate, routage et UTM constituent une base solide pour évoluer vers des fonctions d'ingénieur sécurité réseau.

\section{Mot de fin}

L'expérience vécue chez Axeli a été extrêmement formatrice. Elle m'a permis de concilier les apports théoriques de ma formation à l'EMSI avec les exigences du monde professionnel. Je remercie à nouveau Mr. Othman Lachhab et l'équipe Axeli pour leur confiance et leur accompagnement.
